\begin{table}[t]
\centering
\footnotesize
\caption{Σύνολα πειραμάτων. Κάθε στιγμιότυπο είναι μία διαταραχή ενός γράφου βάσης (Πειράματα Α και Β, Ενότητα~\ref{sec:perturb})· ο διαχωρισμός γίνεται ανά γράφο βάσης. Τα πλήθη στιγμιοτύπων είναι αυτά της επανάληψης του Σταδίου~2 και περιλαμβάνουν τον αδιατάρακτο γράφο A1.0.}
\label{tab:panels}
\begin{tabular}{@{}lllrrr@{}}
\toprule
σύνολο & γεννήτρια & ρόλος & γράφοι & στιγμιότυπα & $n$ \\
\midrule
\texttt{calib} & στατικός ιστός με μπλοκ & προσαρμογή & 64 & 1\,856 & 2\,000 \\
\texttt{test} & στατικός (ανεξάρτητοι σπόροι) & εκτός δείγματος & 32 & 480 & 2\,000 \\
\texttt{scale} & στατικός, $8\times$ μεγαλύτερος & μεταφορά κλίμακας & 12 & 180 & 16\,000 \\
\texttt{topo} & ανάπτυξη / προνομιακή προσκόλληση & μεταφορά τοπολογίας & 24 & 360 & 3\,000 \\
\texttt{real} & Πίνακας~\ref{tab:real} & πραγματικά δίκτυα & 5 & 131 & 2\,995--325\,557 \\
\bottomrule
\end{tabular}
\end{table}
